\chapter{Referencial Teórico}

\section{Data Drift em Sistemas de Aprendizado de Máquina}
Em um cenário onde modelos de aprendizado de máquina são cada vez mais
incorporados a sistemas em produção, a confiabilidade de uma predição é
definida através da estabilidade da relação estatística entre os dados de
treinamento e os dados que o modelo enfrenta em operação. Dessa forma, para
tratar essa questão o estudo de data drift atua como peça fundamental para
compreender por que modelos validados com cuidado se degradam ao longo do
tempo. Segundo \citeonline{gama2014}, o fenômeno consiste na mudança, ao
longo do tempo, da distribuição conjunta entre os dados de entrada e a
variável alvo. Assim, o problema central não é apenas a existência da
mudança, mas o fato de que um modelo continua operando e produzindo saídas
mesmo quando essa distribuição já não representa o presente.

No entanto, o desafio inerente a esse processo reside justamente no caráter
silencioso da degradação. Conforme \citeonline{sculley2015} destacam,
sistemas de aprendizado de máquina incorrem em uma forma de dívida técnica
específica, que se acumula a partir de mudanças no mundo externo e é mais
difícil de detectar do que problemas de código. Esse processo resulta em
perda progressiva de desempenho sem qualquer sinal de erro explícito, pois,
ao contrário de uma falha de código, um modelo degradado por drift permanece
funcional, apenas com qualidade decrescente.

Nesse contexto, é fundamental diferenciar o ambiente de aprendizado em stream
do tradicional aprendizado em lote. Conforme argumentam \citeonline{lu2018},
enquanto o aprendizado em lote pressupõe acesso a um conjunto fixo e completo
de dados antes do treinamento, o aprendizado em stream processa instâncias
que chegam uma a uma, ao longo do tempo. Nesse cenário, o objetivo não é
treinar uma única vez sobre dados estáticos, mas manter o modelo válido
sobre um fluxo potencialmente infinito e sujeito a mudanças.

Diante disso, o desafio central da área reside em decidir como o sistema deve
reagir quando uma mudança é identificada, uma vez que nenhuma estratégia de
resposta se mostra ótima em todos os cenários. Formalmente, adota-se a
definição de \citeonline{gama2014}: dados dois instantes de tempo $t_0$ e
$t_1$, e sendo $p_{t_0}(X, y)$ e $p_{t_1}(X, y)$ as distribuições conjuntas
entre o vetor de atributos $X$ e o rótulo $y$ observadas em cada instante,
diz-se que ocorre drift entre $t_0$ e $t_1$ se existe ao menos um $X$ tal que
$p_{t_0}(X, y) \neq p_{t_1}(X, y)$. Essa formulação, equivalente à adotada
por \citeonline{lu2018}, enquadra o fenômeno em sua definição mais ampla e
sustenta a discussão conceitual que segue: a caracterização dos diferentes
tipos de drift, a lógica das estratégias de retreinamento e a escolha dos
classificadores. O componente empírico deste trabalho, no entanto, observa
apenas o recorte de $p(X)$ acessível ao detector adotado, não a distribuição
conjunta $p(X,y)$ em sua totalidade, distinção retomada com precisão na
Seção~\ref{sec:deteccao}.

\section{Dimensões de Caracterização do Drift}
\label{sec:dimensoes}
Compreendida a definição do fenômeno, o passo seguinte reside em caracterizar
de que forma uma mudança se manifesta, uma vez que essa caracterização é o
que deveria orientar a escolha da resposta. Da literatura emergem quatro
dimensões relativamente independentes, e reconhecê-las é essencial porque um
mesmo episódio de drift pode pertencer simultaneamente a mais de uma delas.

A primeira dimensão é temporal, e descreve como a mudança se desenvolve ao
longo do tempo. Segundo \citeonline{gama2014} e \citeonline{lu2018},
distinguem-se quatro padrões. O drift abrupto ocorre quando a distribuição
muda repentinamente de um estado para outro, dentro de um curto intervalo. O
drift gradual e o incremental descrevem transições mais lentas: no primeiro,
a distribuição alterna entre dois estados até que o novo predomine; no
segundo, os parâmetros mudam em pequenos incrementos contínuos. O drift
recorrente, por sua vez, descreve situações em que uma distribuição
anteriormente observada reaparece após um intervalo, padrão relevante em
domínios com sazonalidade. Conforme demonstram \citeonline{sun2018} e
\citeonline{li2021}, para esse último padrão o reaproveitamento de modelos
históricos é mais eficiente do que o retreinamento do zero.

A segunda dimensão é conceitual, e distingue o que exatamente mudou na
distribuição conjunta entre entradas e rótulo. \citeonline{gama2014}
denominam real drift a mudança na relação entre atributos e rótulo, quando a
própria fronteira de decisão se altera, e virtual drift a mudança restrita à
distribuição das entradas, quando o conceito aprendido permanece válido ainda
que os dados venham de uma região distinta do espaço de atributos. Essa
distinção tem consequência prática direta, pois, conforme demonstram
\citeonline{oliveira2023}, estratégias adequadas para real drift são subótimas
quando aplicadas a virtual drift, precisamente porque, neste último caso, a
relação verdadeira entre entrada e saída não se alterou.

A terceira dimensão é estrutural, e refere-se a onde a mudança se manifesta no
espaço de atributos. Conforme sistematizam \citeonline{barddal2017}, o
conceito de feature drift descreve o fenômeno pelo qual um subconjunto de
atributos deixa de ser relevante, ou passa a ser relevante, para a tarefa de
predição, podendo ocorrer simultaneamente ao drift de conceito tradicional.

A quarta dimensão é quantitativa, e não descreve o que muda, mas o quanto
muda. Conforme propõem \citeonline{wang2024quadcdd}, uma caracterização
completa de um episódio de drift deveria representar conjuntamente seu
início, fim, tipo e severidade, sendo a severidade a magnitude da mudança na
distribuição, ortogonal à sua classificação temporal, conceitual ou
estrutural. \citeonline{wang2022elastic}, no mesmo estudo sobre boosting
adaptativo já citado na introdução, sustentam formalmente essa
relevância ao demonstrar que a severidade do drift determina a taxa de
convergência necessária para a readaptação do modelo, de modo que dois
episódios do mesmo padrão temporal podem exigir respostas de intensidade
distinta caso apresentem severidades diferentes.

A Tabela~\ref{tab:dimensoes} resume as quatro dimensões e os valores que cada
uma pode assumir.

\begin{table}[h]
\centering
\caption{Dimensões de caracterização do drift}
\label{tab:dimensoes}
\begin{tabular}{p{2.6cm}p{3.8cm}p{6.2cm}}
\toprule
\textbf{Dimensão} & \textbf{Valores possíveis} & \textbf{Descrição} \\
\midrule
Temporal & Abrupto, gradual, incremental, recorrente & Como a mudança se
desenvolve ao longo do tempo \\
Conceitual & Real drift, virtual drift & O que muda na distribuição conjunta
entre entrada e rótulo \\
Estrutural & Feature drift & Onde a mudança se manifesta no espaço de
atributos \\
Quantitativa & Severidade (baixa, média, alta) & O quanto a distribuição
muda, independentemente do tipo \\
\bottomrule
\end{tabular}
\end{table}

Diante da convergência dessas quatro dimensões, torna-se evidente por que
caracterizar o tipo de drift é, na prática, mais difícil do que a
formalização sugere. Conforme organizam \citeonline{lu2018}, a pesquisa da
área divide-se em três componentes, a detecção de drift, que responde se a
mudança ocorre, o entendimento de drift, que responde quando, como, onde e
com que intensidade ela ocorre, e a adaptação ao drift, sendo o componente de
entendimento o menos desenvolvido. Esse reconhecimento é central para este
trabalho, cujo objetivo se apoia justamente sobre os sinais que permanecem
observáveis diante dessa dificuldade. Das quatro dimensões apresentadas, este
trabalho manipula experimentalmente a temporal e a quantitativa, únicas
diretamente acessíveis a partir do detector adotado (Seção~\ref{sec:deteccao});
as dimensões conceitual e estrutural permanecem fora do escopo experimental,
tratadas como características latentes não abordadas por este trabalho.

\section{Detecção de Drift}
\label{sec:deteccao}
Estabelecidas as dimensões que caracterizam a mudança, a etapa de
monitoramento é responsável por sinalizar quando ela ocorre. Conforme
sistematiza \citeonline{bayram2022}, os métodos de detecção organizam-se em
duas grandes famílias, distinguidas pela natureza do sinal que monitoram: os
baseados em distribuição, que comparam a distribuição estatística dos dados
de um intervalo de referência com a de um intervalo mais recente, e os
baseados em desempenho, que monitoram a taxa de erro do classificador ao
longo do tempo. A principal vantagem da primeira família é não depender da
disponibilidade imediata de rótulos verdadeiros, exigência que os métodos da
segunda família fazem e que se torna pouco realista em cenários de produção
nos quais o rótulo chega com atraso.

Este trabalho adota o HDDDM (Hellinger Distance Drift Detection Method)
\cite{ditzler2011}, fundamentado em distância de
Hellinger, como detector fixo em todas as condições experimentais, de modo
que a variável observada seja a estratégia de resposta, e não o mecanismo de
detecção.
A escolha entre as três variantes de detectores de distribuição descritas por
\citeonline{bayram2022}: HDDDM, JSDDM (Jensen-Shannon Distance Drift
Detection Method) e KSDDM (Kolmogorov-Smirnov Drift Detection Method),
fundamentadas respectivamente
em distância de Hellinger, divergência de Jensen-Shannon e teste de
Kolmogorov-Smirnov, não é arbitrária: apenas a distância de Hellinger admite
forma fechada geral entre duas distribuições normais, o que permite calcular
a severidade real de cada episódio de drift injetado sem depender de
estimação amostral, questão desenvolvida na Seção~\ref{sec:severidade}.
Comparar um intervalo de referência contra um intervalo recente não é
exclusividade de um cenário de batch offline: o próprio framework genérico de
detecção de \citeonline{lu2018} assume, em seu primeiro estágio, a obtenção de
blocos de dados diretamente a partir de streams, e \citeonline{komorniczak2022}
demonstram viabilidade técnica equivalente para essa família de métodos sob a
denominação explícita de processamento em lote de data streams. Este trabalho
adota, portanto, o mesmo regime de janela deslizante: a cada novo bloco de
instâncias, a divergência entre a janela de referência e a janela mais
recente é recalculada, e o limiar adaptativo original do método é preservado.

Essa escolha, no entanto, tem um custo teórico direto: por operar
exclusivamente sobre $p(X)$, o detector não distingue real drift de virtual
drift \cite{gama2014}, o que mantém a dimensão conceitual do drift como uma
característica latente neste desenho experimental, não estimável a partir do
mecanismo de detecção adotado. O escopo deste trabalho reconhece esse limite
e concentra-se nas dimensões temporal e de severidade, tratadas na Seção a
seguir como as variáveis efetivamente manipuláveis no experimento.

\section{Operacionalização da Severidade}
\label{sec:severidade}
A severidade, para figurar como variável manipulada no experimento fatorial e
não apenas como conceito discutido em abstrato, precisa de uma definição
operacional que permita tanto controlá-la na geração dos dados quanto
compará-la ao que o detector estima a partir da amostra. Este trabalho define
essa operacionalização a partir de um modelo gerador simples e controlável: um
episódio de drift em uma dada feature é modelado como um deslocamento entre
duas distribuições normais de mesma variância, $P = \mathcal{N}(\mu, \sigma^2)$
antes do drift e $Q = \mathcal{N}(\mu + \Delta, \sigma^2)$ depois, onde $\mu$,
$\sigma$ e $\Delta$ são parâmetros escolhidos pelo experimentador.

Sob essa construção, a distância de Hellinger ao quadrado entre $P$ e $Q$,
a mesma métrica sobre a qual o HDDDM opera \cite{ditzler2011}
(Seção~\ref{sec:deteccao}), tem forma fechada geral entre duas normais
univariadas \cite{pardo2006}:
\begin{equation}
H^2(P, Q) = 1 - \sqrt{\frac{2\sigma_P\sigma_Q}{\sigma_P^2+\sigma_Q^2}}
\exp\left(-\frac{(\mu_P-\mu_Q)^2}{4(\sigma_P^2+\sigma_Q^2)}\right).
\end{equation}
Como este trabalho adota variância igual entre $P$ e $Q$, a expressão se
reduz a
\begin{equation}
H^2(P, Q) = 1 - \exp\left(-\frac{\Delta^2}{8\sigma^2}\right),
\label{eq:severidade}
\end{equation}
dependente exclusivamente do deslocamento padronizado $\Delta/\sigma$. A
Equação~\ref{eq:severidade} é invertível: para atingir uma severidade-alvo
$H^2_{\text{alvo}}$, o deslocamento padronizado necessário é
\begin{equation}
\frac{\Delta}{\sigma} = \sqrt{-8\ln\left(1-H^2_{\text{alvo}}\right)}.
\end{equation}
Isso permite definir níveis discretos de severidade a partir de valores-alvo
de $H^2$, por exemplo, severidade baixa ($H^2=0{,}1$, $\Delta/\sigma
\approx 0{,}92$), média ($H^2=0{,}4$, $\Delta/\sigma \approx 2{,}02$) e alta
($H^2=0{,}7$, $\Delta/\sigma \approx 3{,}10$), e de escolher, para cada
condição experimental, o deslocamento exato que produz a severidade
pretendida. A Tabela~\ref{tab:severidade} resume os
três níveis adotados.

\begin{table}[h]
\centering
\caption{Níveis de severidade adotados e o deslocamento padronizado
correspondente}
\label{tab:severidade}
\begin{tabular}{lcc}
\toprule
\textbf{Nível} & \textbf{$H^2_{\text{alvo}}$} & \textbf{$\Delta/\sigma$} \\
\midrule
Baixa & 0,1 & 0,92 \\
Média & 0,4 & 2,02 \\
Alta & 0,7 & 3,10 \\
\bottomrule
\end{tabular}
\end{table}

A severidade assim calculada é a \emph{severidade real}: um valor exato,
obtido diretamente das distribuições geradoras conhecidas, sem qualquer
estimação a partir de amostra. Ela se distingue da \emph{severidade
observada}, que é o valor da mesma métrica de Hellinger calculado pelo HDDDM
sobre a janela de dados amostrados da stream (Seção~\ref{sec:deteccao}),
uma estimativa sujeita ao ruído amostral e ao tamanho de janela adotado. Por
serem exatamente a mesma métrica, aplicada à mesma família de distribuições,
a severidade real e a severidade observada são diretamente comparáveis: a
primeira serve de padrão-ouro para responder à questão principal deste
trabalho (Seção~\ref{sec:questoes}), e a segunda é o que estaria
disponível em produção, onde a distribuição geradora nunca é conhecida;
essa comparação fundamenta a questão derivada.

Vale notar que essa escolha reforça a decisão pelo HDDDM como detector fixo:
o KSDDM admite forma fechada equivalente para a estatística de
Kolmogorov-Smirnov apenas sob a mesma suposição de variância igual, $D_{KS} =
2\Phi(\Delta/2\sigma) - 1$, onde $\Phi$ é a função de distribuição acumulada
da normal padrão, mas o JSDDM não possui forma fechada elementar para a
divergência de Jensen-Shannon entre normais, exigindo integração numérica
mesmo no cenário controlado deste trabalho. O HDDDM é, portanto, o único dos
três cuja severidade real permanece exata sob generalizações futuras do
modelo gerador, como variâncias distintas entre $P$ e $Q$.

\section{Estratégias de Retreinamento}
\label{sec:estrategias}
Uma vez detectada a mudança, resta ao sistema decidir como atualizar o modelo,
etapa que constitui o foco central deste trabalho. Conforme organizam
\citeonline{gama2014}, as estratégias de resposta dividem-se em duas famílias
amplas. A adaptação cega atualiza o modelo continuamente, independentemente de
qualquer sinal explícito de mudança, como em estratégias de janela deslizante
de tamanho fixo. A adaptação informada, por sua vez, condiciona a atualização
à existência de um sinal de detecção, podendo variar desde ajustes
incrementais até a reconstrução completa do modelo, classificada pelos
próprios autores como a reação mais radical possível ao drift.

Para efeitos deste trabalho, adotam-se quatro estratégias operacionais que
instanciam esse espectro, além de uma condição de controle. O reset completo
descarta o modelo corrente e o retreina do zero a partir dos dados mais
recentes, acionado apenas quando o HDDDM sinaliza drift. A janela
deslizante, uma forma de adaptação cega, executa o mesmo procedimento de
retreinamento do zero sobre uma janela de tamanho fixo, mas em intervalos
regulares predefinidos, independentemente de qualquer sinal de detecção;
a diferença entre as duas estratégias está, portanto, inteiramente na
política de acionamento (reativa e esparsa vs. contínua e agendada), não no
procedimento de ajuste do modelo em si, que é idêntico. Como a janela
deslizante retreina em cadência fixa mesmo na ausência de drift, seu custo
computacional acumulado ao longo de uma stream completa tende a ser maior
que o do reset completo, cujo acionamento depende da taxa de detecção do
HDDDM, o oposto do que uma comparação ingênua por evento sugeriria. O
retreinamento incremental ajusta os parâmetros do modelo a cada novo lote de
dados, sem descartar o conhecimento previamente adquirido, também acionado
por detecção. A reativação de modelos históricos, quarta categoria proposta
por \citeonline{sun2018} especificamente para o tratamento de drift
recorrente, mantém um repositório de modelos previamente descartados,
associados cada um à distribuição sob a qual foram treinados, e reativa o
candidato mais próximo da distribuição atual quando um sinal de recorrência
é identificado. Esse sinal é definido pela mesma distância de Hellinger
usada pelo detector: a cada novo bloco, a distribuição atual é comparada
contra a de cada modelo do repositório, e a recorrência é sinalizada quando
a menor distância obtida cai abaixo de um limiar fixo; o modelo
correspondente a essa menor distância é o reativado. Complementa esse
conjunto uma condição de controle sem retreinamento, na qual o modelo
inicial nunca é atualizado independentemente do drift, usada como
referência para avaliar se cada estratégia efetivamente supera a ausência
de adaptação, seguindo o mesmo contraste empregado por \citeonline{rahmani2023}.
A Tabela~\ref{tab:estrategias} compara as cinco condições por três
critérios operacionais.

\begin{table}[h]
\centering
\caption{Comparação operacional entre as estratégias de retreinamento e a
condição de controle}
\label{tab:estrategias}
\begin{tabular}{p{2.8cm}p{2.6cm}p{2.8cm}p{2.8cm}}
\toprule
\textbf{Estratégia} & \textbf{Acionada por detecção?} & \textbf{Descarta
conhecimento anterior?} & \textbf{Custo computacional acumulado} \\
\midrule
Reset completo & Sim & Sim & Baixo a médio (depende da taxa de detecção) \\
Janela deslizante & Não (cadência fixa) & Sim, a cada atualização & Alto
(retreina mesmo sem drift) \\
Retreinamento incremental & Sim & Não & Baixo \\
Reativação de modelos históricos & Sim, com sinal de recorrência & Não
(reaproveita modelo) & Baixo \\
Sem retreinamento (controle) & Não se aplica & Não se aplica & Nulo \\
\bottomrule
\end{tabular}
\end{table}

Essas quatro estratégias, contrastadas com a condição de controle,
raramente são comparadas sob as mesmas condições
experimentais, o que motiva a comparação sistemática proposta neste trabalho,
sob condições controladas de tipo, severidade e classificador.

No entanto, a eficácia de cada estratégia não é uniforme. Conforme demonstram
\citeonline{rahmani2023} com dados reais, o desempenho do retreinamento depende
do tipo de drift enfrentado e da disponibilidade de rótulos para sustentá-lo,
com o retreinamento adequado recuperando e superando o desempenho de um modelo
sem atualização na maioria dos cenários testados. Diante disso, reforça-se a
premissa central deste trabalho, de que a estratégia de resposta ao drift não
deveria ser tratada como uma escolha única e fixa, mas condicionada às
características da mudança e do modelo em uso.

\section{Classificadores}
Além da estratégia de resposta, o próprio classificador monitorado influencia
a forma como o sistema reage ao drift. Conforme demonstram
\citeonline{wang2022dynamic}, em estudo distinto sobre ponderação dinâmica
de atributos em streams (não se confunde com o estudo sobre boosting
adaptativo citado nas seções anteriores), um mesmo mecanismo de adaptação
produz ganhos de
desempenho distintos para diferentes classificadores, o que motiva a
investigação dessa interação a partir de modelos com mecânicas distintas.

O Naive Bayes é um classificador probabilístico fundamentado no teorema de
Bayes, sob a suposição de independência condicional entre os atributos dado o
rótulo. Essa suposição torna o algoritmo leve e adequado a ambientes de
stream, mas também potencialmente mais sensível a mudanças nas correlações
entre atributos, já que o modelo assume, por construção, que essa correlação
não existe.

O Random Forest é um método de conjunto que combina múltiplas árvores de
decisão, cada uma treinada sobre uma amostra distinta dos dados e sobre um
subconjunto aleatório de atributos em cada divisão de nó. Conforme discutem
\citeonline{gomes2017}, variantes adaptativas desse método para aprendizado em
stream incorporam mecanismos de detecção de drift internamente ao próprio
conjunto de árvores, substituindo árvores obsoletas conforme o drift é
identificado.

O XGBoost (eXtreme Gradient Boosting), por fim, é um método de conjunto baseado em boosting, no qual cada
nova árvore é treinada para corrigir os erros cometidos pelo conjunto já
construído. Originalmente concebido para aprendizado em lote, sua inclusão em
ambiente de stream requer uma adaptação que constrói novos membros do conjunto
a partir de pequenos lotes de dados conforme estes se tornam disponíveis.

\section{Sinais Observáveis e Resposta ao Drift}
\label{sec:sinais}
Reunidos os elementos anteriores, torna-se possível articular o problema que
motiva este trabalho. Conforme discutido, o tipo de drift que deveria orientar
a escolha da estratégia é, em produção, uma característica latente: o sistema
não observa diretamente se a mudança é abrupta ou gradual, real ou virtual.
Diante disso, o que se observa são sinais indiretos, aproximações imperfeitas
das características verdadeiras da mudança.

Alguns desses sinais guardam relação razoavelmente direta com as dimensões do
drift. A velocidade com que a estatística de um detector muda entre janelas
sucessivas é um indício da dimensão temporal, enquanto a semelhança entre a
distribuição atual e distribuições históricas, definida operacionalmente na
Seção~\ref{sec:estrategias} como a menor distância de Hellinger entre a
distribuição atual e o repositório de modelos descartados, indica
recorrência. A magnitude
da própria estatística de divergência calculada pelo HDDDM, por sua vez, é a
severidade observada definida na Seção~\ref{sec:severidade}, um indício de
severidade já demonstrado em detectores de distribuição por
\citeonline{komorniczak2022}, cuja relevância como variável de resposta é
sustentada formalmente pelo mesmo estudo sobre boosting adaptativo de
\citeonline{wang2022elastic} discutido anteriormente. A comparação entre
essa severidade observada e a severidade real, calculada por fórmula fechada
(Seção~\ref{sec:severidade}), é o que sustenta o protocolo exploratório
descrito na metodologia.

A dimensão conceitual, no entanto, mostra-se mais resistente a essa tradução.
Conforme reconhecem \citeonline{bayram2022}, o único sinal razoavelmente
direto para distinguir uma mudança real de uma mudança apenas na distribuição
das entradas é a taxa de erro do classificador, disponível somente quando
rótulos verdadeiros chegam em tempo hábil. Como o detector adotado neste
trabalho pertence à família baseada em distribuição (Seção~\ref{sec:deteccao}),
esse sinal de erro não integra o mecanismo de detecção, o que mantém a
dimensão conceitual como uma característica latente, fora do escopo deste
trabalho. Diante desse limite, a orientação preliminar buscada como
desdobramento exploratório concentra-se nas dimensões temporal e de
severidade, únicas para as quais os sinais observáveis do detector adotado
oferecem uma aproximação razoável, aspecto que fundamenta o desenho
experimental descrito na metodologia deste trabalho.